\section{The complete local two-branch object and the original finite conductor}
\label{sec:real-pair-local-algebra}

This is a derivation for the certified member of the original geometric
family, with its original parameters recovered by exact inverse maps.
The point and full-series bounds are RPZ1--19 \cite{RPZ005};
the conjugation, old four-label return and finite reciprocal calculation
are RPCJ1--23, RPE1--26 and RPAC1--15
\cite{RPCJ005,RPE005,RPAC005}.
The companion complete directional source, RPD1--31 including RPD12a
\cite{RPD006,RPD31Source},
supplies the exact nonzero real Jacobian used below.  Its identity
RPD4 also follows directly from RPZ10--15, as recalled here.
The original matrix and Gamma frames are WCF1--6 \cite{WCF}, whose
\href{https://github.com/KokunoYumeto/zeta-function-research-reader/blob/5992ad50c0f2a06921f1fcaded7b4e38097c0739/workbenches/splitzero-tandem/continuations/20260920-original-kernel-mixed-determinants/providers/WEIGHTED_CONDUCTOR_FORWARD.tex#L28}{complete pinned proof}
retains the Meixner--Pollaczek provenance, DLMF
(18.19.8)--(18.19.9), (18.22.8), (18.23.7).
The original factor family is PCL1--10 \cite{PCLOriginal}, with its
\href{https://github.com/KokunoYumeto/zeta-function-research-reader/blob/487253b5b01851da5b9b98af4e146fecf1fab902/workbenches/splitzero-tandem/continuations/20260919-native-conductor-coprimality/providers/PCL_COMPLETE.tex#L19}{complete pinned proof}.
We prove the local constructions and all finite module identities below;
elementary analytic inversion, polynomial division and matrix elimination
are used with their proofs, not claimed as new general theorems.
The four marked states retain the canonical reader by The Clankers,
its theorem \texttt{explicit-four-time-Jacobian-collision}, and its
distinction between the antecedent three-dimensional example and the
four-dimensional marked construction \cite{ES488}.  The certified
interval arithmetic used in RPZ is Fredrik Johansson's Arb
\cite{human:arb2016}; the classical Gamma dictionary is
\cite{DLMF18Source}.  No arithmetic identification is made.

\subsection{Original functions, coefficient reflection and the exact parameter inverse}

Let $p=(z,x)$ and let $p_*=(z_*,x_*)$ be the unique real root in the
exact rational rectangle RPZ4.  Retain $u=1/z$, the four primary
exponents
\[
 \begin{gathered} (\rho_{\mathrm M}^{\rm prim},\rho_{\mathrm N}^{\rm prim},
 \rho_{\mathrm T}^{\rm prim},\rho_{\mathrm E}^{\rm prim})
 =(3/4+3i,3/4-3i,1/4+3i,1/4-3i),\\
 \delta=1/4,\quad\gamma=3,\quad\beta_0=143/8,\quad\eta=21025/256.
 \end{gathered}\tag{RLA1}
\]
The subscript on $\beta_0$ distinguishes this unchanged original
quartic constant from the spectral root $\beta(p)$ constructed below.
Likewise $\alpha_{\rm unit}$ denotes the unchanged nonzero real
scale written $\alpha$ in RPZ3: the actual unit remains
$a=\alpha_{\rm unit}(1+ix)$.
Put $w_a=\rho_a^{\rm prim}-1/2$, $V_{ar}=w_a^r$,
$J_d=\operatorname{diag}(1,-1,-1,1)$, $W(x)=I+ixJ_d$,
$\zeta=e^{2\pi i/5}$ and $D_j=\operatorname{diag}(\zeta^{-ja})_{a=1}^4$.
The complete original functions, with $1\le a\le4$, $0\le r\le3$, are
\[
\begin{aligned}
 \mathcal R(z)&=\sum_{n\ge0}R_nz^n,\\
 (R_n)_{ar}&=\sum_{\substack{h,k\ge0\\2h+4k=5n+r+1-a}}
 \frac{(\beta_0/3)^h\eta^k(-5)^{h+k-n}(a/5)_{h+k-n}}{h!k!},\\
 M_j(p)&=W(x)^{-1}V\mathcal R(z)^{-1}D_j\mathcal R(z)V^{-1}W(x),\\
 f_j(\xi,p)&=Q_0\bigl(M_j(p)(e^{\rho_a^{\rm prim}\xi})_a\bigr),\\
 Q_0(y)&=y_1y_4-y_2y_3,\\
 G(\xi,p)&=f_1(\xi,p)f_2(\xi,p),\\
 G^\#(\xi,p)&=\overline{G(\bar\xi,\bar p)},\\
 E_A(\xi,p)&=G(\xi,p)G^\#(\xi,p),\\
 g(t,p)&=e^{-4\omega(t)}E_A(\omega(t),p),\\
 \omega(t)&=-i\arctan t.
\end{aligned}\tag{RLA2}
\]
Empty sums are zero.  RPZ5--9 prove convergence, holomorphicity and
matrix invertibility on a complex neighborhood of $p_*$ after shrinking
it.  This last assertion also follows directly from the entire
convergence bound and nonzero determinants at $p_*$.
The reflection in RLA2 conjugates the Taylor coefficients in all three
variables.  It is holomorphic; it is not pointwise conjugation at a
complex parameter.  The real-coefficient series $\mathcal R$,
$\bar V=J_cV$, $\overline{W(\bar x)}=J_cW(x)J_c$, and
$Q_0(J_cy)=-Q_0(y)$ give $f_{5-j}=-f_j^\#$ identically.
Here $J_c$ exchanges indices $1,2$ and $3,4$.
Thus $G^\#=f_4f_3$ and the two minus signs cancel in $E_A$.
This proves the factor identity throughout the complex neighborhood.

All following derivatives without displayed arguments are at $(0,p_*)$.
Set $b_*=f_2\ne0$, $B=G_\xi\ne0$, $A_z=G_z$, $A_x=G_x$,
$b=|B|^2>0$ and
\[
 \begin{gathered}\Delta=\Im(\overline{A_z}A_x)\ne0,\qquad
 \mathcal H(p)=\binom{G(0,p)}{G^\#(0,p)},\\
 J_{\mathbb C}=D\mathcal H(p_*)=
 \begin{pmatrix}A_z&A_x\\\overline{A_z}&\overline{A_x}\end{pmatrix},
 \quad\det J_{\mathbb C}=-2i\Delta.
 \end{gathered}\tag{RLA3}
\]
The nonvanishing statement is proved, rather than assumed: RPZ12--13
give $\|I-PJ\|<1$ for the full real Jacobian
$J=D_p(\Re f_1,\Im f_1)$ and RPZ11 gives $\det P\ne0$.
The convergent geometric matrix series in $I-PJ$ inverts $PJ$,
so $J$ is invertible.  At the zero, multiplication by $b_*$ takes
$J$ to $D_p(\Re G,\Im G)$ and has determinant $|b_*|^2$.
Its determinant is precisely $\Delta$, proving RLA3.
The quantitative enclosure RPD31 has $\Delta<0$, but nonvanishing
alone suffices for the construction.

There is an actual holomorphic inverse $p=\mathcal P(r_+,r_-)$ to
$\mathcal H$, with
\[
 J_{\mathbb C}^{-1}
 =\frac1{-2i\Delta}
 \begin{pmatrix}\overline{A_x}&-A_x\\-\overline{A_z}&A_z\end{pmatrix},
 \qquad
 p_{n+1}=p_n-J_{\mathbb C}^{-1}
              \bigl(\mathcal H(p_n)-(r_+,r_-)^{\mathsf T}\bigr),
 \quad p_0=p_*.
 \tag{RLA4}
\]
To prove it, choose a closed complex ball about $p_*$ on which
$\|I-J_{\mathbb C}^{-1}D\mathcal H\|\le1/2$; continuity and RLA3
give such a ball.  Restrict $(r_+,r_-)$ so the displacement of its
center under RLA4 is at most one quarter of that ball's radius.
Integration along a segment gives a contraction ratio $1/2$ and
an image inside three quarters of the radius.  The differences of
successive iterates are bounded by a geometric series.  The uniform
limit exists and is the unique fixed point, hence exactly a solution
of $\mathcal H(p)=(r_+,r_-)$.
Every iterate is holomorphic in the parameters; local uniform
convergence and the coefficientwise Cauchy integral formula show that
the limit is holomorphic.  Uniqueness proves both inverse compositions
after restriction to these neighborhoods.  In particular all original
coordinates are recovered, with $u=1/\mathcal P_z$ and
$a=\alpha_{\rm unit}(1+i\mathcal P_x)$, without identifying a scale
or an original parameter with a chosen constant.

Since $g(0,p)=r_+r_-$, its zero divisor is the union of the two
smooth curves $r_+=0$ and $r_-=0$, meeting transversely only at $p_*$.
Each factor has multiplicity one.  On the real parameter slice
$r_-=\overline{r_+}$, so their union meets that slice only at $p_*$.
This is the complete local zero divisor of the original zeroth
moment, not only its tangent cone.

\subsection{The two actual spectral roots and their full analytic units}

Put $F_+(t,p)=G(\omega(t),p)$ and
$F_-(t,p)=G^\#(\omega(t),p)$.
Their $t$ derivatives at the base point are $-iB$ and
$-i\bar B$, respectively, both nonzero.  There are unique
holomorphic roots $\alpha(p),\beta(p)$ close to zero, given by the
convergent iterations
\[
\begin{aligned}
 t^+_{n+1}&=t^+_n-F_+(t^+_n,p)/(-iB),&t^+_0&=0,\\
 t^-_{n+1}&=t^-_n-F_-(t^-_n,p)/(-i\bar B),&t^-_0&=0,\\
 \alpha(p)&=\lim_nt^+_n,&\beta(p)&=\lim_nt^-_n.
\end{aligned}\tag{RLA5}
\]
The proof is the one-dimensional version of RLA4: on a sufficiently
small fixed complex disk and parameter neighborhood each derivative
of the iteration in its first variable has modulus at most $1/2$,
and its value at zero has modulus at most one quarter of the radius.
The same geometric and Cauchy integral arguments give unique
holomorphic limits.  Consequently these are roots of the actual
infinite-series factors.  Because $f_2$ and $f_3$ stay nonzero, they
are exactly the roots of the original $f_1$ and $f_4$, respectively.
The original spectral coordinates are
$\xi_+=\omega(\alpha)$ and $\xi_-=\omega(\beta)$; the inverse
spectral morphism is $t=\tan(i\xi)=i\tanh\xi$ on this neighborhood.
Multiplication of its derivative by $\omega'(t)$ and the values at
zero, or the defining arctangent identity, verifies both compositions.

The complete units and complete factorization are
\[
\begin{aligned}
 U_+(t,p)&=\int_0^1\partial_tF_+
       (\alpha(p)+\theta(t-\alpha(p)),p)\,d\theta,\\
 U_-(t,p)&=\int_0^1\partial_tF_-
       (\beta(p)+\theta(t-\beta(p)),p)\,d\theta,\\
 F_+&=(t-\alpha)U_+,&F_-&=(t-\beta)U_-,\\
 U(t,p)&=e^{-4\omega(t)}U_+(t,p)U_-(t,p),&
 q(t,p)&=(t-\alpha)(t-\beta),\\
 g(t,p)&=U(t,p)q(t,p),&U(0,p_*)&=-|B|^2=-b.
\end{aligned}\tag{RLA6}
\]
Integration of the derivative proves both divisions exactly.
Integration of locally convergent Taylor series proves the units
holomorphic; their values at the base point are $-iB$ and
$-i\bar B$.  They, their product, and the center exponential are
nonzero on a smaller neighborhood.  No coefficient of any unit
has been removed.  In particular
\[
 r_+=-\alpha U_+(0,p),\qquad r_-=-\beta U_-(0,p),\qquad
 g_0=U(0,p)\alpha\beta.
 \tag{RLA7}
\]
Thus $\alpha=0$ and $\beta=0$ are precisely the two original
parameter branches, with the original $r_\pm$ factors retained.

The roots themselves are a pair of holomorphic coordinates.  Implicit
differentiation of the two proved zero equations gives
\[
 D\alpha(p_*)=-\frac{i}{B}(A_z,A_x),\qquad
 D\beta(p_*)=-\frac{i}{\bar B}(\bar A_z,\bar A_x),\qquad
 \det D(\alpha,\beta)(p_*)=\frac{2i\Delta}{b}\ne0.
 \tag{RLA8}
\]
An exact inverse, not just the derivative, is constructed as follows.
For proposed roots $(A_0,B_0)$ close to zero iterate
\[
 p_{n+1}=p_n-J_{\mathbb C}^{-1}
       \binom{F_+(A_0,p_n)}{F_-(B_0,p_n)},\qquad p_0=p_*.
 \tag{RLA9}
\]
Its parameter derivative vanishes at $(A_0,B_0,p)=(0,0,p_*)$.
The ball argument of RLA4 proves a unique holomorphic limit
$p=\mathcal P_{\rm roots}(A_0,B_0)$.
Uniqueness in RLA5 then proves
\[ (\alpha,\beta)(\mathcal P_{\rm roots}(A_0,B_0))=(A_0,B_0). \]
The other inverse composition follows from uniqueness in RLA9.
This explicitly restores $z,x,u,a$ from the spectral-root coordinates.

For a parameter function write $h^\sharp(p)=\overline{h(\bar p)}$.
Since $\overline{\omega(-\bar t)}=\omega(t)$, the exact reflection
and real-slice restrictions are
\[
 F_-(t,p)=\overline{F_+(-\bar t,\bar p)},\qquad
 \beta=-\alpha^\sharp,\qquad
 U(-\bar t,\bar p)=\overline{U(t,p)},\qquad
 \beta(p)=-\overline{\alpha(p)}\quad(p\in\mathbb R^2).
 \tag{RLA10}
\]
The first identity and uniqueness in RLA5 prove the second.
Reflecting both linear factors in RLA6, including their minus signs,
proves the identity for $U$.  On the real slice $U(0,p)$ is real
and stays negative, since its value at $p_*$ is $-b$.
Accordingly $g_0=-U(0,p)|\alpha(p)|^2>0$ away from $p_*$.
The spectral collision curve $\alpha=\beta$ is a different curve
from the parameter branches $\alpha=0$, $\beta=0$.
At its nonzero points $g_0\ne0$, so collision of the two nonzero
spectral roots does not make the finite conductor singular.

\subsection{Every cutoff and the full two-generator module}

Let $R$ be the local ring of convergent holomorphic parameter germs
at $p_*$.  RLA9 gives mutually inverse ring identifications
$R\simeq\mathbb C\{\alpha,\beta\}$; the original parameter maps and
unit remain those already displayed.  Put $\sigma=\alpha+\beta$,
$\pi=\alpha\beta$, $R_N=R[t]/(t^N)$ for $N\ge1$ and
$q=t^2-\sigma t+\pi$.  Write $L_N(h)$ for multiplication by the
full series $h$ modulo $t^N$ in the ordered basis $1,t,\ldots,t^{N-1}$.
Each entry uses the exact coefficient, not a truncated substitute
for that coefficient.  Since $U(0,p)$ is a unit, the finite inverse
of $L_N(U)$ is multiplication by $U^{-1}$ modulo $t^N$; its
coefficients are specified recursively by
\[
 \begin{gathered}v_0=U_0^{-1},\qquad
 v_n=-U_0^{-1}\sum_{j=1}^nU_jv_{n-j},\\
 L_N(g)=L_N(q)L_N(U),\qquad
 \mathcal M_N:=\operatorname{coker}L_N(g)=R[t]/(t^N,q).
 \end{gathered}\tag{RLA11}
\]
Here $U_j=[t^j]U$, and the equality of cokernels uses the identity on
the target: $(g)=(q)$ in $R_N$, because multiplication by $U$ is
invertible there.  The corresponding source morphism is $L_N(U)$
and its exact inverse is the displayed recurrence.  Its metric
effect is calculated below, rather than declared isometric.

Define $h_{-1}=0$, $h_0=1$ and
\[
 h_n(\alpha,\beta)=\sum_{j=0}^n\alpha^{n-j}\beta^j\quad(n\ge0),
 \qquad h_n=\sigma h_{n-1}-\pi h_{n-2}\quad(n\ge1),\qquad
 C=\begin{pmatrix}0&-\pi\\1&\sigma\end{pmatrix}.
 \tag{RLA12}
\]
The recurrence for $n=1$ means $h_1=\sigma$; for $n\ge2$ it
follows by multiplying and cancelling the finite sums.
Polynomial division by the monic $q$ is exact over $R$ and leaves
a unique representative $a+bt$.  Indeed each power $t^n$ with
$n\ge2$ is reduced using $t^2=\sigma t-\pi$, and a nonzero
multiple of the monic quadratic has degree at least two.
Multiplication by $t$ on this rank-two module has the matrix $C$.
Induction gives the complete remainders
\[
 [t^n]_q=-\pi h_{n-2}+h_{n-1}t\quad(n\ge1),\qquad
 C^N=\begin{pmatrix}
 -\pi h_{N-2}&-\pi h_{N-1}\\h_{N-1}&h_N
 \end{pmatrix}\quad(N\ge1).
 \tag{RLA13}
\]
The first formula at $n=1$ uses $h_{-1}=0$.
The columns of $C^N$ are exactly the remainders of $t^N,t^{N+1}$.
Every multiple of $t^N$ in $R[t]/q$ is their $R$-linear combination,
since that ring has basis $1,t$.  Hence the full conductor has the
two-generator presentation
\[
 R^2\xrightarrow{\ C^N\ }R^2\longrightarrow\mathcal M_N
 \longrightarrow0,\qquad
 [a+bt]\longmapsto[a+bt]\bmod(t^N,q).
 \tag{RLA14}
\]
The inverse to its cokernel map reduces any polynomial modulo $q$
by RLA13.  Reducing a multiple of $t^N$ changes its remainder by
$C^N$ times a vector, and reducing a multiple of $q$ changes it by
zero.  These facts prove well-definedness and both compositions.
Conversely, a remainder in the image of $C^N$ is a multiple of
$t^N$ modulo $q$, so the same argument proves injectivity.
Thus RLA14 describes the entire module, not only its dimension at
the special point.  It includes $N=1$: the unit entry of $C$ leaves
$\mathcal M_1=R/(\pi)$.

In the displayed presentation the determinant and entry ideals are
\[
 \operatorname{Fitt}_0(\mathcal M_N)=(\pi^N),\qquad
 \operatorname{Fitt}_1(\mathcal M_N)
   =(h_{N-1},\pi h_{N-2}),\qquad
 \operatorname{Fitt}_j(\mathcal M_N)=R\quad(j\ge2).
 \tag{RLA15}
\]
Here Fitting ideals mean the ideals of minors of sizes two, one
and zero in RLA14.  The determinant is $(\det C)^N=\pi^N$;
$h_N=\sigma h_{N-1}-\pi h_{N-2}$ gives the entry ideal.
For $N=1$ it is $R$, since $h_0=1$.
These are also the corresponding ideals from the original
$N$-generator presentation: adding a generator with a relation that
solves uniquely for it, or eliminating such a pair, adjoins or removes
an identity block.  The minors of that enlarged block give precisely
the same ideals with their shifted sizes.  Multiplication of a
presentation by invertible matrices leaves its minors' ideals fixed
by the determinant expansion, applied also to the inverse matrices.
For an explicit common presentation start with generators
$1,t,\ldots,t^{N+1}$ and relations
$q,tq,\ldots,t^{N-1}q,t^N,t^{N+1}$.
Eliminating the last two generators with their unit-coefficient
relations gives the original multiplication-by-$q$ presentation
modulo $t^N$.  Instead eliminating $t^2,\ldots,t^{N+1}$ in order
with the first $N$ monic relations leaves $1,t$ and precisely the
two columns of $C^N$.  These are the indicated invertible
generator and relation operations; no ideal is inferred from
dimensions alone.
For $N\ge2$ the common zero of the two generators of
$\operatorname{Fitt}_1$ near the origin is just $\alpha=\beta=0$:
if $\pi=0$ and one root is nonzero, $h_{N-1}$ is its nonzero
$(N-1)$st power; if $\pi\ne0$, consecutive $h$'s cannot both vanish,
as the recurrence backwards would force $h_0=0$.

\subsection{The exact gluing of the two branches}

Evaluation at the two roots is an injective map on $R[t]/q$:
\[
 a+bt\longmapsto(a+b\alpha,a+b\beta),\qquad
 \mathcal L=\{(v_+,v_-)\in R^2:
                      v_+-v_-\in(\alpha-\beta)\}.
 \tag{RLA16}
\]
Its image is exactly $\mathcal L$: divide $v_+-v_-$ by
$\alpha-\beta$ to recover $b$, and then set $a=v_+-b\alpha$.
Its kernel is zero since $(\alpha-\beta)b=0$ in the integral
domain of convergent series.  Multiplication by $t^N$ becomes
$\operatorname{diag}(\alpha^N,\beta^N)$ on this lattice.
It follows, with all maps specified, that
\[
\begin{gathered}
 0\longrightarrow\mathcal M_N
 \xrightarrow{\,\iota\,}R/(\alpha^N)\oplus R/(\beta^N)
 \xrightarrow{\,d\,}R/(\alpha-\beta,\alpha^N)
 \longrightarrow0,\\
 \iota([P])=(P(\alpha)\bmod\alpha^N,
                         P(\beta)\bmod\beta^N),\qquad
 d([v_+],[v_-])=[v_+-v_-].
\end{gathered}\tag{RLA17}
\]
The last map is well defined and onto.  Its kernel consists of
pairs whose representatives differ by an element of
$(\alpha-\beta,\alpha^N)$; changing the first representative by
a multiple of $\alpha^N$ puts them in $\mathcal L$, proving that
the kernel is the image of $\iota$.
For injectivity, suppose $(\alpha^Nr,\beta^Ns)\in\mathcal L$.
Modulo $\alpha-\beta$ the relation is
$\alpha^N(r-s)=0$ in $\mathbb C\{\alpha\}$.
This ring is a domain, so $r-s$ is divisible by $\alpha-\beta$.
Thus $(r,s)\in\mathcal L$, and the original pair is a multiple of
$t^N$ in the lattice, proving that its class is zero.
This proves exactness directly, including the intersection coupling.
The gluing quotient has complex dimension $N$, with basis
$1,\alpha,\ldots,\alpha^{N-1}$ after setting $\beta=\alpha$.
The module is not the direct sum of the two branch modules for
any $N\ge1$.  For $N=1$ its minimal generator count is one,
whereas that direct sum needs two, as reduction modulo
$(\alpha,\beta)$ shows.  For $N\ge2$, RLA15's first Fitting
ideal contains $h_{N-1}$, a nonzero homogeneous polynomial of
degree $N-1$.  The first Fitting ideal of the diagonal
presentation of $R/(\alpha^N)\oplus R/(\beta^N)$ is
$(\alpha^N,\beta^N)$; every element has order at least $N$.
Those ideals differ, so the presentation invariance just proved
excludes an isomorphism.  RLA17 retains the exact embedding and
gluing map that relates these modules.

The annihilator is exactly
\[
 \operatorname{Ann}_R(\mathcal M_N)
       =(\alpha^N)\cap(\beta^N)=(\alpha^N\beta^N).
 \tag{RLA18}
\]
Injection in RLA17 gives inclusion from right to left.
Each projection of $\iota$ is onto, because constant polynomials
give every residue in each component; hence an annihilator must
lie in both ideals.  In the convergent double Taylor series,
divisibility by $\alpha^N$ removes all powers of $\alpha$ below
$N$, and divisibility by $\beta^N$ removes all powers of $\beta$
below $N$.  Both together give divisibility by their product.
At a point of $\alpha=0$ with $\beta\ne0$, the second branch
module and the gluing quotient in RLA17 vanish locally, so the
module is exactly $R/(\alpha^N)$ there.  The corresponding
statement on the other branch is $R/(\beta^N)$.
At their intersection, RLA14 reduced modulo $(\alpha,\beta)$ has
zero matrix for $N\ge2$, and therefore a two-dimensional fiber.
For $N=1$ the fiber has dimension one, as already calculated.
The two dimensions at the intersection record the two original
lost directions; the gluing map retains how the branches approach
that intersection.

\subsection{Exact connection to the directional resonances and the full inverse}

Away from $\alpha\beta=0$ the exact inverse coefficients are
\[
 [t^n]\frac1{g(t,p)}
 =\sum_{j=0}^n v_j(p)
       \frac{h_{n-j}(\alpha(p),\beta(p))}
            {(\alpha(p)\beta(p))^{n-j+1}},\qquad 0\le n<N,
 \tag{RLA19}
\]
where every $v_j$ is the full unit reciprocal from RLA11.
Indeed the geometric expansions of $(t-\alpha)^{-1}$ and
$(t-\beta)^{-1}$ yield
$[t^m]q^{-1}=h_m(\alpha,\beta)/(\alpha\beta)^{m+1}$.
Finite convolution with $U^{-1}$ proves RLA19.  It holds also
when $\alpha=\beta\ne0$, where the polynomial finite sum gives
$h_m=(m+1)\alpha^m$; no divided-difference singularity is inserted.
This formula includes every higher unit coefficient and the complete
$e^{-4\omega}$ center factor in RLA6.

For a nonzero real vector $v$ in the unchanged $(z,x)$ coordinates,
let $A_v=A_zv_z+A_xv_x$ and define the literal derivatives
$a_v=-iA_v/B$, $b_v=-i\bar A_v/\bar B=-\bar a_v$.
They are nonzero because the real derivative is invertible.  Along
$p_*+\tau v$, RLA8 gives
$\alpha=\tau a_v+O(\tau^2)$ and
$\beta=\tau b_v+O(\tau^2)$, while $v_0(p_*)=-1/b$.
Consequently RLA19 proves the same exact leading reciprocal
coefficient as RPD13--15:
\[
 H_D(v)=-\frac{h_D(a_v,b_v)}{b(a_vb_v)^{D+1}},\qquad
 H_D(v)=0\ \Longleftrightarrow\
 \frac{a_v}{b_v}=e^{2\pi ik/(D+1)},\quad1\le k\le D.
 \tag{RLA20}
\]
For the first equality, the $j=0$ term in RLA19 has order
$\tau^{-D-2}$; each $j>0$ term has strictly lower order.
For the second, $a_v,b_v\ne0$ and
$(a_v-b_v)h_D(a_v,b_v)=a_v^{D+1}-b_v^{D+1}$.
At equality $a_v=b_v$, $h_D=(D+1)a_v^D\ne0$; otherwise the
displayed roots of unity are all possibilities.
Since $a_v/b_v=(A_v/B)/\overline{(A_v/B)}$, this recovers exactly
the original resonance lines of RPD20, including their orientation
and original scale.  Thus the old reciprocal resonance is the
vanishing of the degree-$D$ remainder coefficient of the actual
two-root algebra, and the other entry of the presentation remains.
RLA19, rather than just its leading term, retains the curvature
and unit corrections calculated in RPD23--30.

\subsection{Original Gamma metric and an exact residual matrix using every coefficient}

For $s\in\{1,a_{\rm amp}\}$, $c_0=a_{\rm amp}/2$, $c_1=c_0-4$,
retain the original data
\[
\begin{gathered}
 dm_s(y)=\frac{(2\pi)^{s/2}2^{s/2}}{4\pi\Gamma(s/2)}
       |\Gamma(s/4+iy/2)|^2dy,\quad M_s=(2\pi)^{s/2},\quad
 \rho_n=\sqrt{n!/(s/2)_n},\\
 p_{0,s}=1,\quad p_{1,s}=y,\quad
 p_{n+1,s}=yp_{n,s}-n(n-1+s/2)p_{n-1,s},\\
 h_{n,s}=M_sn!(s/2)_n,\qquad
 \phi_{n,s,c}(S)=p_{n,s}((S-c)/i)/\sqrt{h_{n,s}},\qquad
 \|P\|_{s,c}^2=\int|P(c+iy)|^2dm_s(y).
\end{gathered}\tag{RLA21}
\]
These are WCF2's exact original orthonormal frames, with no mass
or center replacement.  For $D=N-1$ set $D_\rho=\operatorname{diag}
(\rho_0,\ldots,\rho_D)$ and let $J_N$ reverse these $N$ indices.
If $x$ is the source orthonormal coefficient column, the exact
polynomial coordinate morphism and inverse are
\[
 x\longmapsto a(t)=\sum_{j=0}^{N-1}(J_ND_\rho x)_jt^j,
 \qquad x=D_\rho^{-1}J_N\operatorname{coeff}(a),\qquad
 B_{D,s}=D_\rho^{-1}J_NL_N(g)J_ND_\rho.
 \tag{RLA22}
\]
The same formula uses the target frame at $c_1$ on the output.
Since $L_N(g)_{ij}=g_{i-j}$ for $i\ge j$, reversal makes its
entries $g_{j-i}$ above the diagonal, proving RLA22 entry by entry.
The original Gamma norm in these polynomial coordinates is
\[
 a^*W_{N,s}a,\qquad
 W_{N,s}=\operatorname{diag}(\rho_{N-1}^{-2},\ldots,\rho_0^{-2}).
 \tag{RLA23}
\]
Here $a$ in the quadratic form means its coefficient column;
the original unit still has the distinct meaning specified in RLA1.
Formula RLA23 is the exact transport of the original normalized
frame coefficients, not an assertion that the physical measure
has lost its mass.  Its inverse in RLA22 uses the full $h_{n,s}$
and centers in RLA21 when it reconstructs the physical polynomial.
For the source change $c=L_N(U)a$ in RLA11, the source Gram becomes
\[
 (L_N(U)^{-1})^*W_{N,s}L_N(U)^{-1},
 \tag{RLA24}
\]
whereas the target Gram remains $W_{N,s}$.
This explicitly accounts for every unit coefficient that affects
the cutoff; no multiplication by a nonconstant unit is asserted
to be an isometry.

There is also a two-by-two residual presentation directly in the
original orthonormal matrix, retaining all the higher $g_j$.
For $N\ge2$ order the columns as $(2,\ldots,N-1,0,1)$ and keep
the row order $(0,\ldots,N-1)$.  With the corresponding permutation
matrix $\Pi$, write
\[
 B_{D,s}\Pi=\begin{pmatrix}P&Q\\R_0&S_0\end{pmatrix},\qquad
 P_{ij}=B_{i,j+2}\quad(0\le i,j\le N-3),\qquad
 \Sigma_N=S_0-R_0P^{-1}Q.
 \tag{RLA25}
\]
The remaining blocks are the literal remaining rows and columns,
so this formula specifies every entry using RLA22.
At the base point $P$ is upper triangular with nonzero diagonal
$-b\rho_{i+2}/\rho_i$.  Thus $\det P$ is a unit germ and
$P^{-1}=\operatorname{adj}(P)/\det P$ is an exact finite formula.
For $N=2$ the pivot block is empty and $\Sigma_2=B_{1,s}$.
The two invertible elimination matrices satisfy
\[
 E_L=\begin{pmatrix}I&0\\-R_0P^{-1}&I_2\end{pmatrix},\qquad
 E_R=\begin{pmatrix}P^{-1}&-P^{-1}Q\\0&I_2\end{pmatrix},\qquad
 E_L(B_{D,s}\Pi)E_R=\operatorname{diag}(I_{N-2},\Sigma_N).
 \tag{RLA26}
\]
Block multiplication proves this identity and its inverse; all
entries, including $g_3,\ldots,g_{N-1}$, enter $P,Q,R_0,S_0$.
It proves directly $\operatorname{coker}\Sigma_N\simeq\mathcal M_N$.
Moreover the column permutation has sign $(-1)^{2(N-2)}=1$, so
\[
 \det\Sigma_N=\frac{g_0^N}{\det P}
       =\frac{U(0,p)^N(\alpha\beta)^N}{\det P},\qquad
 G_{\rm source}=E_R^*E_R,\qquad
 G_{\rm target}=(E_L^{-1})^*E_L^{-1}.
 \tag{RLA27}
\]
Indeed original source coordinates are $\Pi E_R$ times the new
ones and original target coordinates are $E_L^{-1}$ times the
new ones; the original frames are orthonormal.
These complete Grams include all cross terms between the pivot
coordinates and the two residual coordinates.
For example every squared original singular value is exactly a
root in $\lambda$ of
\[
 \det\left(
 \operatorname{diag}(I,\Sigma_N)^*G_{\rm target}
 \operatorname{diag}(I,\Sigma_N)-\lambda G_{\rm source}\right)=0.
 \tag{RLA28}
\]
This is the defining eigenvalue equation for the transported
Hermitian form and follows by multiplying by an invertible square
root of $G_{\rm source}$.  Thus the two-generator algebra and the
original analytic norm are connected by exact finite morphisms;
the residual matrix alone is never substituted for the full metric.
At $N=1$ the original matrix is the scalar $g_0$, already covered
by RLA7 and RLA11--18.

\subsection{All three conductor domains and the four original ES labels}

The exact order $v=\operatorname{ord}_{t=0}g(t,p)$ is determined
on the whole local complex parameter space by
\[
\begin{array}{c|c|c}
 \text{parameter locus}&v&g_v\\\hline
 \alpha\beta\ne0&0&U_0\alpha\beta\\
 \alpha=0,\ \beta\ne0&1&-U_0\beta\\
 \beta=0,\ \alpha\ne0&1&-U_0\alpha\\
 \alpha=\beta=0&2&U_0=-b
\end{array}.
 \tag{RLA29}
\]
All entries follow by multiplying the complete factorization RLA6.
Since $\omega'(0)=-i\ne0$, this is also the exact order of
$E_A(\xi,p)$ at $\xi=0$, with $g_v=(-i)^v\mu_v/v!$.
For every $D\ge0$ the actual order-adjusted map is therefore
\[
 \mathcal A_{D,s,p}:\mathcal P_{D+v}/\mathcal P_{v-1}
       \longrightarrow\mathcal P_D,\qquad
 (\mathcal A_{D,s,p})_{ij}=
 \begin{cases}g_{j-i+v}(p)\rho_{j+v}/\rho_i,&j\ge i,\\0,&j<i.
 \end{cases}
 \tag{RLA30}
\]
Its diagonal $g_v\rho_{j+v}/\rho_j$ never vanishes on its stated
stratum.  The source quotient uses the original norm at $c_0$ and
the target uses the original norm at $c_1$.  WCF3's generating
identity, or its coefficient comparison in WCF4--6, proves this
matrix; inversion is the exact weighted triangular reciprocal of
$t^{-v}g$.  The fixed old map on $\mathcal P_D$, in contrast, has
kernel $\mathcal P_{\min(v,D+1)-1}$ and image
$\mathcal P_{D-v}$ when $D\ge v$, and zero image when $D<v$.
This follows immediately from
$T_AS^n=\sum_j\binom nj\mu_jS'^{n-j}$ and nonzero leading
$\binom nv\mu_v$ for $n\ge v$.
The two parameter branches have $v=1$ individually, their
intersection has $v=2$, and the real slice has only $v=0,2$ locally.
Both changed domains and the old kernel remain explicit.

Retain the four original ES columns $K$, ordered
$\mathrm M,\mathrm N,\mathrm T,\mathrm E$, and the exact Lagrange
polynomials
$\ell_a(S)=\prod_{b\ne a}(S-\rho_b^{\rm prim})/
(\rho_a^{\rm prim}-\rho_b^{\rm prim})$.
Let $\Phi_0q=\sum_a(K^{-1}q)_a\ell_a(S)$; thus
$\Phi_0^{-1}P=K(P(\rho_a^{\rm prim}))_a$.
On the stratum of order $v=0,1,2$, the exact augmented old map is
\[
 \Theta_vq=\left(T_A\Phi_0q,
          ([S^r]\Phi_0q)_{0\le r<v}\right)
   \in\mathcal P_{3-v}\oplus\mathbb C^v.
 \tag{RLA31}
\]
It is an isomorphism.  To give its inverse for every stratum,
write the supplied target as $Y=\sum_{k=0}^{3-v}Y_kS'^k$ and
set $c_r=t_r$ for $0\le r<v$.  Solve successively, in descending
order $k=3-v,\ldots,0$, by the exact finite equations
\[
 c_{k+v}=\frac{Y_k-
       \sum_{n=k+v+1}^3\binom nk\mu_{n-k}c_n}
                   {\binom{k+v}{k}\mu_v},\qquad
 \Theta_v^{-1}(Y,t)=K\left(\sum_{n=0}^3c_n
                                    (\rho_a^{\rm prim})^n\right)_a.
 \tag{RLA32}
\]
All denominators are nonzero by RLA29.
The coefficient of $S'^k$ in $T_A\sum c_nS^n$ is the numerator
equation used to solve for $c_{k+v}$; the lower coefficients
$c_0,\ldots,c_{v-1}$ contribute zero.  This proves the two inverse
compositions by descending substitution and interpolation.
At $v=2$, it is exactly the RPE20--21 map, with the two original
lost scalars $t_0,t_1$.  At $v=1$ it reads explicitly
\[
 c_3=\frac{Y_2}{3\mu_1},\quad
 c_2=\frac{Y_1-3\mu_2c_3}{2\mu_1},\quad
 c_1=\frac{Y_0-\mu_2c_2-\mu_3c_3}{\mu_1},\quad c_0=t_0.
 \tag{RLA33}
\]
Thus each individual branch loses one scalar and the intersection
loses the two-dimensional old kernel, without merging any ES label.
For every $v$ the four complete states are
\[
 \Theta_v(Ke_a)=
 \left(T_A\ell_a,([S^r]\ell_a)_{r<v}\right),\qquad
 x_a^{(v)}=[S^v\ell_a]\in\mathcal P_{v+3}/\mathcal P_{v-1},\qquad
 y_a^{(v)}=T_A(S^v\ell_a).
 \tag{RLA34}
\]
The four $x_a^{(v)}$ are independent because division by $S^v$
on the representative with degrees $v$ through $v+3$ recovers
the independent Lagrange basis.  RLA30 then proves independence
of the four corrected outputs.  With $X_ve_a=x_a^{(v)}$,
$Y_ve_a=y_a^{(v)}$, the exact morphisms remain
$T_AX_vK^{-1}=Y_vK^{-1}$, and conjugation of the unchanged ES
polynomial through these invertible maps gives the same original
intertwining as RPE18.  RLA31--33 separately recover the old
fixed source with its lost scalars.

For completeness every original metric cross term in this recovery
is fixed by a finite formula.  Let $S_v$ be the coefficient matrix
of RLA32 from $(Y_0,\ldots,Y_{3-v},t_0,\ldots,t_{v-1})$ to
$(c_0,c_1,c_2,c_3)$.  Put $m_0=1$, $m_2=s/2$,
$m_4=(s/2)(3s/2+2)$, $m_6=(s/2)(15(s/2)^2+30(s/2)+16)$,
and $m_1=m_3=m_5=0$.  Then the complete transported source Gram is
\[
 G_{\Theta_v}=S_v^*H_{s,c_0}S_v,\qquad
 (H_{s,c})_{jk}=M_s\sum_{a=0}^j\sum_{d=0}^k
 \binom ja\binom kd c^{j+k-a-d}(-i)^ai^d m_{a+d},
 \quad0\le j,k\le3.
 \tag{RLA35}
\]
The finite sum expands
$\int\overline{(c+iy)^j}(c+iy)^kdm_s(y)$.
The six moments are those obtained from WCF2's original recurrence
and norms, as calculated in RPE22--23; substituting them gives
every entry, including the source mass and all cross terms.
The output metric uses the corresponding upper-left
$(4-v)\times(4-v)$ block of $H_{s,c_1}$ instead.
Thus neither the original target metric nor either lost scalar is
silently replaced by a quotient coordinate metric.

Equations RLA2--35 construct one complete local object: the actual
original analytic family with its invertible parameter maps, its
two transverse parameter branches, its full analytic spectral units,
and every finite weighted conductor and branch-glued cokernel.
All roots, branches and domains belong to this certified geometric
member.  No equation identifies it with an arithmetic zeta quartet.
